% 本文件由 main.tex 输入，作为“引言”章节，不单独编译。

\section{引言}

创新是推动高质量发展的重要动力，科技创新能力提升已经成为现代化产业体系建设中的关键内容。近几年，围绕科技创新、产业升级和新质生产力培育，政府投资基金和创业投资在政策层面受到持续重视。《国务院办公厅关于印发〈促进创业投资高质量发展的若干政策措施〉的通知》提出，要围绕创业投资“募投管退”全链条完善政策环境，强化企业创新主体地位，引导创业投资更好发挥投早、投小、投长期、投硬科技的作用。随后，《国务院办公厅关于促进政府投资基金高质量发展的指导意见》进一步明确，政府投资基金要突出政府引导和政策性定位，按照市场化、法治化、专业化原则规范运作，发展耐心资本，更好服务现代化产业体系建设，同时还强调要合理统筹基金布局，防止同质化竞争和对社会资本产生挤出效应。在这样的政策背景下，政府产业引导基金已经不只是一般意义上的财政工具，而是被赋予了支持科技创新、撬动社会资本和优化产业结构的重要功能。

从现实发展来看，政府产业引导基金扩张较快，已经成为地方政府支持创新和招商引资的重要手段。已有研究表明，政府引导基金可以通过提高失败容忍度、缓解融资约束和集聚社会资本等方式促进企业创新，尤其是在关键核心技术突破、区域创新投入增长等方面表现出一定的积极作用（吴超鹏和严泽浩，2023；蔡庆丰等，2024）。不过，政府产业引导基金的政策目标能否稳定实现，并不是一个没有条件的问题。与政策文本中对“耐心资本”“投早投小”和“市场化运作”的强调相比，地方政府在现实中还面临着财政收支矛盾加大、债务风险防控趋严和预算约束增强等压力。特别是《国务院办公厅关于促进政府投资基金高质量发展的指导意见》已经明确提出，政府投资基金布局和运行要统筹考虑财力基础、产业基础以及债务风险，这说明地方债务约束并不是外在背景，而是已经进入政府投资基金高质量发展的政策框架之中。

也正因为如此，一个值得进一步讨论的问题就变得比较清楚了：当地方政府同时承担稳增长、促创新和防风险等多重任务时，债务压力会不会改变政府产业引导基金的运作逻辑，进而影响其对创新的促进作用。按政策设计看，政府产业引导基金应当通过股权投资方式支持高风险、长周期的创新活动，并通过让利和增信机制撬动更多社会资本进入科技创新领域；但在债务压力较大的情形下，地方政府可能更重视基金安全性、回收速度和短期绩效，这样一来，基金原本应有的耐心资本属性、社会资本撬动功能以及融资约束缓解作用，就有可能被削弱。换句话说，政府产业引导基金是否促进创新，不能只从基金本身来判断，还需要把它放到地方财政债务约束不断增强的现实情境中来考察。

现有研究已经分别从两个方面积累了较多成果。一方面，关于政府产业引导基金与创新的研究表明，引导基金既可能通过融资支持、信号传递和资本集聚等机制促进创新，也可能因为代理问题、风险规避和目标偏离而削弱创新激励。另一方面，关于财政压力和地方政府债务的研究则指出，财政收支约束和债务扩张会通过挤出金融资源、提高融资成本、强化短期激励和扭曲政府干预方式等渠道抑制企业创新与长期投资。相比之下，把“地方债务压力—政府产业引导基金运作—创新绩效”放在同一分析框架中的研究还比较少。现有文献更多是在讨论基金是否有效，或者财政压力是否抑制创新，而对债务压力如何改变政府基金的行为边界，特别是如何影响其耐心资本属性、社会资本撬动效率和融资约束缓解功能，仍缺少较为系统的分析。

基于此，本文将地方债务压力、政府产业引导基金与创新绩效纳入统一分析框架，重点讨论两个层面的问题：第一，政府产业引导基金总体上能否促进创新；第二，在地方债务压力上升的背景下，这种促进作用是否会被削弱，以及这种削弱主要通过哪些机制实现。围绕上述问题，本文进一步从耐心资本属性、社会资本撬动效率和融资约束缓解效应三个方面展开分析，力图更完整地说明政府产业引导基金为什么在某些地区和条件下能够较好促进创新，而在另一些地区和条件下其效果又会明显减弱。

本文的研究可能有以下几方面意义。第一，现有研究多把政府产业引导基金视为独立的创新政策工具，本文则进一步把地方债务压力纳入分析，有助于从财政约束角度理解基金创新效应的差异来源。第二，已有研究虽然讨论了财政压力抑制创新的问题，但对政府产业引导基金这一中观政策工具的传导作用关注不多，本文有助于拓展地方债务压力影响创新的机制链条。第三，从现实政策角度看，两份国务院文件一方面强调发展耐心资本、支持科技创新，另一方面也强调统筹债务风险、防止同质化竞争和挤出效应，本文的讨论能够在一定程度上回应这种政策上的内在张力，并为理解政府产业引导基金何以“有效”、又何以可能“失效”提供补充证据。

下文其余部分安排如下：第二部分为文献综述与理论假说，第三部分为研究设计，第四部分为实证结果分析，第五部分为结论及进一步讨论。